\chapter{Gallbladder Pain}

\section*{Definition}
Gallbladder pain is usually right-upper abdominal or upper-midline pain caused by gallstones, inflammation, or impaired bile drainage. It may follow a fatty meal and radiate to the back or right shoulder.

\section*{When to See a Doctor}
Seek urgent evaluation for fever, jaundice, chills, persistent severe pain, repeated vomiting, confusion, fainting, or low blood pressure. Arrange prompt assessment for persistent, recurrent, unexplained, or function-limiting symptoms.

\section*{How It Develops}
The symptom reflects irritation, inflammation, obstruction, altered secretion, infection, or dysfunction in the relevant organ system. The pattern, duration, associated symptoms, exposures, and examination findings determine whether it is self-limited or requires investigation.

\section*{Causes}
\begin{itemize}
  \item Biliary colic, acute cholecystitis, bile-duct stone, cholangitis, pancreatitis, and gallbladder malignancy.
  \item Medication effects, environmental exposures, and chronic disease may contribute.
  \item Less common but important causes include malignancy, vascular disease, or organ failure when symptoms persist or progress.
\end{itemize}

\section*{Related Symptoms}
\begin{itemize}
  \item Pain, fever, fatigue, nausea, or changes in normal organ function
  \item Bleeding, weight change, shortness of breath, weakness, or neurologic symptoms when a serious cause is present
\end{itemize}

\section*{Medical Evaluation}
\begin{itemize}
  \item History addresses onset, duration, severity, triggers, exposures, medications, medical conditions, pregnancy status when relevant, and warning symptoms.
  \item Examination focuses on vital signs and the organ systems suggested by the symptom.
  \item Testing may include blood or urine studies, cultures, imaging, physiologic testing, or specialist examination according to the clinical pattern.
\end{itemize}

\section*{Treatment Options}
Treatment depends on the cause: uncomplicated biliary colic, cholecystitis, cholangitis, and pancreatitis require different management. Persistent or severe pain may need urgent imaging, antibiotics, drainage, or surgery under specialist care.

\section*{Self-Care and Home Management}
Avoid known triggers, maintain reasonable hydration, record timing and associated symptoms, and use only treatments appropriate for the suspected cause. Do not use leftover antibiotics or delay urgent care because symptoms temporarily improve.

\section*{References}
\begin{thebibliography}{99}
\bibitem{gallbladder_pain:cholecystitis} Trowbridge RL, Rutkowski NK, Shojania KG. Does this patient have acute cholecystitis? \textit{JAMA}. 2003;289:80--86.
\bibitem{gallbladder_pain:tokyo} Yokoe M, Hata J, Takada T, et al. Tokyo Guidelines 2018: diagnostic criteria and severity grading of acute cholecystitis. \textit{J Hepatobiliary Pancreat Sci}. 2018;25:41--54.
\bibitem{gallbladder_pain:easl} European Association for the Study of the Liver. EASL Clinical Practice Guidelines on the prevention, diagnosis and treatment of gallstones. \textit{J Hepatol}. 2016;65:146--181.
\end{thebibliography}
